\documentclass[a4paper,10pt]{exam}
\usepackage{fontspec}
\setmainfont{OpenDyslexic}[Scale=0.9]  % solo testo; XeLaTeX o LuaLaTeX
\usepackage[italian]{babel}

% --- TIPOGRAFIA ELEGANTE: PALATINO E MATEMATICA ARMONIZZATA ---
\usepackage{newpxmath}
\usepackage{microtype}

\usepackage{amsfonts}
\usepackage{amsmath}
\usepackage{graphicx}
\usepackage{siunitx}
\usepackage{physics}
\usepackage{geometry}
\usepackage{enumitem}
\usepackage[table]{xcolor}
\usepackage{booktabs}
\usepackage{tcolorbox}
\usepackage{tabularx}
\usepackage{multirow}
\usepackage{tikz}

% Risoluzione compatibilità siunitx e physics
\AtBeginDocument{\RenewCommandCopy\qty\SI}
\newcommand{\um}[2]{\SI[output-decimal-marker={,}]{#1}{#2}}

\bracketedpoints
\pointpoints{punto}{punti}
\marginpointname{ punti}
\boxedpoints

%%%%%%%%%%%%%%%%%%%%%%%%%%%%%%%%%%%%%%%%%%%%%%%%%%%%%%%%%%%%%
%%%%%%%%%%%%%%%%%% GEOMETRIA PAGINA %%%%%%%%%%%%%%%%%%%%%%%%%
%%%%%%%%%%%%%%%%%%%%%%%%%%%%%%%%%%%%%%%%%%%%%%%%%%%%%%%%%%%%%
\geometry{
   top=1.0cm,
   bottom=1.0cm,
   left=0.8cm,
   right=0.8cm
}

\pagestyle{headandfoot}
% INTESTAZIONE VUOTA: Divieto assoluto di titolo, materia, indicazione fila
\firstpageheader{}{}{}
\firstpagefooter{}{\scriptsize Pagina \thepage\ di \numpages}{}
\runningheader{}{}{}
\runningfooter{}{\scriptsize Pagina \thepage\ di \numpages}{}

%%%%%%%%%%%%%%%%%%%%%%%%%%%%%%%%%%%%%%%%%%%%%%%%%%%%%%%%%%%%%
%%%%%%%%%%%%%%%%%% STILE MONOCROMATICO %%%%%%%%%%%%%%%%%%%%%%
%%%%%%%%%%%%%%%%%%%%%%%%%%%%%%%%%%%%%%%%%%%%%%%%%%%%%%%%%%%%%
% REGOLA: Solo toni di bianco, nero e sfumature di grigio (gray!*, black!*)

% Badge compatto per numerazione degli esercizi
\newcommand{\badgeQ}[1]{%
  \tikz[baseline=(char.base)]{%
    \node[shape=rectangle, rounded corners=2pt, fill=gray!15, draw=gray!40, line width=0.4pt, inner sep=2.5pt, minimum width=1.6em] (char) {\small\textbf{#1}};%
  }%
}
\renewcommand{\questionlabel}{\badgeQ{\thequestion}\enspace}
\renewcommand{\partlabel}{\textbf{\alph{partno})}}

% Caselle di spunta (checkbox) per griglie e quesiti
\newcommand{\cbox}{\ensuremath{\vcenter{\hbox{\tikz{\draw[black!75, line width=0.5pt, rounded corners=1.2pt] (0,0) rectangle (3.0mm, 3.0mm);}}}}}

% Box consegne/istruzioni con sfondo tenue senza intestazione "Istruzioni"
\tcbset{
  istruzionibox/.style={
    colback=gray!6,
    colframe=gray!30,
    arc=2pt,
    boxrule=0.5pt,
    left=8pt, right=8pt, top=5pt, bottom=5pt,
    fontupper=\small
  }
}

%%%%%%%%%%%%%%%%%%%%%%%%%%%%%%%%%%%%%%%%%%%%%%%%%%%%%%%%%%%%%
%%%%%%%%%%%%%%%%%% GRIGLIE DI VALUTAZIONE %%%%%%%%%%%%%%%%%%%
%%%%%%%%%%%%%%%%%%%%%%%%%%%%%%%%%%%%%%%%%%%%%%%%%%%%%%%%%%%%%

% --- OPZIONE 1: Griglia Competenze con Legenda a Destra ---
\newcommand{\grigliacompetenze}{
\vspace{1mm}
\noindent
\begin{minipage}[c]{0.62\linewidth}
\footnotesize
\setlength{\aboverulesep}{0pt}
\setlength{\belowrulesep}{0pt}
\renewcommand{\arraystretch}{1.25}
\begin{tabular}{l cc >{\columncolor{gray!20}}c ccc c}
\toprule
\rowcolor{gray!10}
\textbf{Indicatori} & \textbf{N} & \textbf{Q} & \cellcolor{gray!20}\textbf{S} & \textbf{D} & \textbf{B} & \textbf{O} & \textbf{Voto} \\
\midrule
Analizzare    & \cbox & \cbox & \cbox & \cbox & \cbox & \cbox & \multirow{4}{*}{\rule{1.4cm}{0pt}} \\
Sviluppare    & \cbox & \cbox & \cbox & \cbox & \cbox & \cbox & \\
Argomentare   & \cbox & \cbox & \cbox & \cbox & \cbox & \cbox & \\
Comunicare    & \cbox & \cbox & \cbox & \cbox & \cbox & \cbox & \\
\bottomrule
\end{tabular}
\end{minipage}\hfill
\begin{minipage}[c]{0.35\linewidth}
\scriptsize \color{black!75}
\textbf{Legenda giudizi:}\par\vspace{1mm}
\begin{tabular}{@{}ll@{}}
\textbf{N}: Non sufficiente & \textbf{D}: Discreto \\
\textbf{Q}: Quasi sufficiente & \textbf{B}: Buono \\
\textbf{S}: Sufficiente & \textbf{O}: Ottimo \\
\end{tabular}
\end{minipage}
}

% --- OPZIONE 2: Griglia Recupero per Argomenti (senza legenda) ---
\newcommand{\linearecupero}[1]{#1 & \cbox & \cbox & \cbox & \\}

\newcommand{\grigliarecupero}[1]{
\vspace{1mm}
\noindent
\begin{center}
\footnotesize
\renewcommand{\arraystretch}{1.2}
\begin{tabular}{p{0.54\textwidth} ccc p{0.24\textwidth}}
\toprule
\rowcolor{gray!12}
\textbf{Nucleo tematico / Obiettivo} & \textbf{NR} & \textbf{PR} & \textbf{R} & \textbf{Note / Esito} \\
\midrule
#1
\bottomrule
\end{tabular}
\end{center}
}

% Contenuto facoltativo: soltanto costanti, formule o supporti approvati.
\newcommand{\supportiprova}[1]{%
  \par\noindent\begin{minipage}{\linewidth}\small
  #1
  \end{minipage}\par\vspace{3mm}%
}

%%%%%%%%%%%%%%%%%%%%%%%%%%%%%%%%%%%%%%%%%%%%%%%%%%%%%%%%%%%%%
\begin{document}
% intestazione sempre presente e non modificabile salvo esplicita richiesta dell'utente
\begin{tcolorbox}[istruzionibox]
È consentito l'uso di penne nere o blu (non è consentito l'uso del lapis e del bianchetto).
È consentito l'uso della calcolatrice scientifica non grafica.
Scrivere i passaggi in modo chiaro e leggibile, esplicitando formule, passaggi fondamentali e motivando le scelte principali.
\end{tcolorbox}

\vspace{1.5mm}


\begin{questions}
\question Nel piano cartesiano sono dati i punti $A(-2,1)$, $B(4,3)$ e $C(2,-3)$.
\begin{parts}
\part Scrivi l'equazione della retta $AB$ e calcola la lunghezza del segmento $AB$.
\part Calcola la distanza di $C$ dalla retta $AB$ e l'area del triangolo $ABC$.
\part Scrivi l'equazione della retta per $C$ parallela ad $AB$ e determina in quale punto interseca l'asse $x$.
\end{parts}
\question Date le rette $r:\ y=3x+2$ e $s:\ kx+(k+1)y-5=0$, determina il valore di $k$ per cui $s$ è perpendicolare a $r$. Discuti il caso in cui $s$ non ha coefficiente angolare.
\question Scrivi l'equazione dell'asse del segmento di estremi $P(-1,4)$ e $Q(5,0)$ e verifica che un suo punto, a tua scelta, è equidistante da $P$ e da $Q$.
\question Dato il fascio di rette $k\,(2x+y+1)+(-x+2y-3)=0$, determina il centro del fascio, la retta del fascio che passa per il punto $(3,0)$ e quella parallela all'asse $x$.
\end{questions}
\null
\vfill

\supportiprova{\textbf{Formulario}\\[1mm]
\begin{tabular}{@{}ll@{}}
Distanza fra due punti & $\sqrt{(x_2-x_1)^2+(y_2-y_1)^2}$ \\
Punto medio & $\left(\dfrac{x_1+x_2}{2},\dfrac{y_1+y_2}{2}\right)$ \\
Distanza punto-retta & $\dfrac{|ax_0+by_0+c|}{\sqrt{a^2+b^2}}$ con $ax+by+c=0$ \\
Rette parallele e perpendicolari & $m_1=m_2$\quad $m_1\,m_2=-1$ \\
Fascio di rette & $k\,(a_1x+b_1y+c_1)+(a_2x+b_2y+c_2)=0$ \\
\end{tabular}}
\noindent
\textbf{Cognome e Nome:}\enspace\hrulefill\hspace{1cm}%
\textbf{Classe:}\enspace\makebox[2.5cm]{\hrulefill}\hspace{1cm}%
\textbf{Data:}\enspace\makebox[2.5cm]{\hrulefill}

\vspace{3mm}
\grigliacompetenze
\end{document}
